% \chapter{Conclusões e Trabalhos Futuros}
% \label{chap:conclusoes-e-trabalhos-futuros}

% Parte final do texto na qual se apresentam as conclusões apoiadas no desenvolvimento do assunto. É a recapitulação sintética dos resultados obtidos. Pode apresentar recomendações e sugestões para pesquisas futuras.

% %\label{sec:contribuicoes-do-trabalho}



% %\label{sec:limitacoes}

% ==============================================================================
% CONCLUSÕES E TRABALHOS FUTUROS
% Números conforme resultados_redes.tex e resultados_marcadores.tex.
% Revisar o intervalo das redes (0,48--0,56) quando as CNNs de 2 min terminarem.
% ==============================================================================

\chapter{Conclusões e Trabalhos Futuros}
\label{cap:conclusoes}

Esta dissertação investigou se o registro de Holter de 24~horas contém
informação capaz de identificar, entre pacientes com insuficiência cardíaca
em ritmo sinusal, aqueles que sofrerão morte súbita cardíaca em quatro anos.
Sobre os 517 pacientes da coorte MUSIC, com 55 casos, foram comparadas três
abordagens no mesmo protocolo de validação cruzada por paciente: redes
neurais treinadas diretamente sobre o sinal das três derivações, modelos
sobre marcadores eletrocardiográficos extraídos do Holter e modelos sobre
variáveis clínicas. Por fim, avaliou-se uma quarta via, em que a rede aprende
primeiro a prever atributos clínicos disponíveis para todos os pacientes e
sua representação é depois usada na predição da morte súbita.

\section{Conclusões}
\label{sec:conc-conclusoes}

A resposta obtida, dentro das condições deste estudo, é negativa para o
sinal bruto. Nenhuma das redes neurais avaliadas --- redes convolucionais
de três capacidades, redes residuais de 18, 34 e 50 camadas, uma rede com
mecanismo de atenção e redes configuradas por busca de hiperparâmetros ---
discriminou casos e controles de forma distinguível do acaso. As \gls{AUC}
fora da partição ficaram entre 0,48 e 0,56, sem relação com a capacidade da
rede, com a duração da janela de entrada ou com a arquitetura, e a variação
entre duas execuções da mesma rede foi da mesma ordem que a variação entre
arquiteturas. Os registros de treinamento mostraram, em todos os casos, a
memorização dos pacientes de treinamento nas primeiras épocas.

Os marcadores eletrocardiográficos tiveram desempenho modesto. Sobre os
44 resumos dos segmentos de 5~minutos, os modelos alcançaram \gls{AUC}
entre 0,585 e 0,624, com intervalos que excluem o acaso por pequena margem
e abaixo do modelo clínico. Quando poucos marcadores foram escolhidos dentro
de cada partição, a \gls{AUC} foi de 0,580; escolhidos sobre a coorte
inteira, alcançaram valores maiores, mas otimistas. Já dois marcadores
definidos antes da análise, pela literatura --- a variação da frequência
cardíaca ao longo do dia e a carga de ectopia ---, alcançaram 0,647 sobre os
segmentos de 30~segundos, sem esse otimismo. Os marcadores de segmentos de 30~segundos
mostraram, em grande parte, diferenças significativas apenas quando cada
segmento foi tratado como observação independente, efeito da
pseudorreplicação, e desapareceram quando a análise foi feita por paciente.

As variáveis clínicas foram a fonte de informação mais consistente, com
\gls{AUC} de 0,661 [0,577--0,741], e a adição dos marcadores ou do sinal
não trouxe ganho significativo. A rede de fusão confirmou esse resultado de
forma direta: com o extrator de sinal congelado, ela obteve o mesmo
desempenho que com o extrator treinado com taxa reduzida, e o desempenho
caiu quando o extrator foi treinado livremente.

A estratégia de rótulos substitutos mudou o resultado para o sinal. Treinada
para prever 45 atributos do ecocardiograma, do \gls{ECG} de repouso, do
Holter e do laboratório, uma ResNet-18 produziu uma representação que, sem
nenhuma variável clínica, alcançou \gls{AUC} de 0,682, equivalente à das
variáveis clínicas, contra 0,508 da rede treinada diretamente para o
desfecho. Essa representação, contudo, não acrescentou informação aos exames
medidos (diferença de $-0{,}005$), e os atributos de maior valor
prognóstico --- tamanho do átrio esquerdo, Pro-BNP e \gls{FEVE} --- foram os
que a rede aprendeu pior. Sem nenhum rótulo, o aprendizado autossupervisionado
encontrou apenas padrões fracos (\gls{AUC} de 0,590 no método contrastivo).
O melhor modelo do trabalho combinou as variáveis clínicas com os exames do
estudo, em especial o ecocardiograma (\gls{AUC} de 0,733).

Dois resultados complementares ajudam a interpretar o quadro. Quando a
restrição ao ritmo sinusal foi removida, a mesma rede residual passou a
discriminar acima do acaso, o que mostra que o pipeline aprende quando há
um padrão reconhecível no sinal --- no caso, a fibrilação atrial --- e
justifica a restrição adotada. E, na literatura, o desempenho de um modelo
de aprendizado profundo pré-treinado em grandes bases caiu para 0,60 no
subgrupo de insuficiência cardíaca \cite{kim2025deep}, valor próximo dos
obtidos aqui.

A conclusão central é, portanto, que, com 33 casos de treinamento por
partição, não foi possível aprender do sinal do Holter um padrão preditivo
de morte súbita quando a rede foi treinada diretamente para o desfecho; que o
sinal contém informação prognóstica comparável à das variáveis clínicas,
acessível quando o aprendizado é orientado por atributos clínicos, mas que
essa informação coincide com a dos exames já disponíveis; e que as estimativas otimistas descritas
na literatura e observadas em etapas iniciais deste trabalho decorrem, em
boa parte, de escolhas de avaliação: seleção de variáveis sobre a coorte
inteira, divisões únicas com poucos casos no teste, tratamento de segmentos
como observações independentes e inclusão de ritmos não sinusais. A
principal contribuição da dissertação está em quantificar esses efeitos
sobre os mesmos pacientes e em documentar, com intervalos de confiança, os
resultados de todas as configurações avaliadas, e não apenas os da melhor.

\section{Limitações}
\label{sec:conc-limitacoes}

O número de eventos é a principal limitação: 55 casos no total e 11 por
conjunto de teste tornam os intervalos de confiança largos, de modo que
diferenças de até 0,1 na \gls{AUC} não podem ser detectadas com segurança.
A coorte é única, sem validação externa, e cada paciente tem um único
registro de Holter. A busca de hiperparâmetros não usou validação cruzada
aninhada, e as redes residuais foram treinadas com uma configuração
anterior à das redes convolucionais. Por fim, o desfecho foi tratado como
classificação binária em quatro anos, o que descarta a informação sobre o
tempo até o evento e exclui os pacientes vivos com seguimento mais curto.
A rede multitarefa foi treinada com uma única semente, e os pacientes de
treinamento da segunda etapa foram vistos pela rede na primeira, ainda que
sem o rótulo de morte súbita.

\section{Trabalhos futuros}
\label{sec:conc-futuros}

Os resultados apontam para direções em que o problema do número reduzido de
eventos é contornado, e não apenas enfrentado com arquiteturas maiores:

\begin{enumerate}
    \item \textbf{Representações pré-treinadas.} Treinar o extrator de
          características em bases grandes de \gls{ECG}, por aprendizado
          auto-supervisionado ou com rótulos substitutos abundantes, e ajustar sobre
          a coorte MUSIC apenas um classificador pequeno, como em
          \citeonline{kim2025deep}.

    \item \textbf{Foco nos atributos prognósticos.} Na rede multitarefa,
          dar maior peso ao tamanho do átrio esquerdo, ao Pro-BNP e à
          \gls{FEVE}, e usar janelas mais longas ou batimentos médios, que
          preservam melhor a onda P, ligada ao aumento do átrio esquerdo.

    \item \textbf{Marcadores da repolarização.} Incluir os marcadores que
          \citeonline{ramirez2017sudden} identificaram como os mais informativos
          nesta coorte --- alternância da onda T e restituição da morfologia da
          onda T ---, que exigem delineação precisa e não foram calculados aqui.

    \item \textbf{Análise de sobrevivência.} Modelar o tempo até o evento,
          com riscos competitivos entre morte súbita e morte por falência de bomba,
          o que aproveita todos os pacientes e todo o seguimento.

    \item \textbf{Validação externa.} Avaliar os modelos em outra coorte de
          pacientes com insuficiência cardíaca e Holter de 24~horas, condição
          necessária para qualquer conclusão sobre generalização.
\end{enumerate}





